\begin{abstract}
Dataset distillation seeks to compress large-scale datasets into compact synthetic surrogates that preserve training utility. Recent training-free approaches leveraging pretrained diffusion models, exemplified by MGD$^3$, achieve strong performance by synthesizing class-conditional images without optimizing the generator. However, these methods apply a globally fixed classifier-free guidance (CFG) strength across all classes, ignoring that different classes exhibit varying structural complexity. We propose \cags{} (Class-Adaptive Guidance Strength), a training-free method that replaces the globally fixed guidance with per-class adaptive guidance based on a multi-factor complexity metric. Through systematic factor ablation, we find that per-class adaptation---not merely the presence of guidance---drives the improvement: cluster entropy is the most effective single complexity factor, while inter-class separability is the least effective, producing near-constant guidance across classes. We provide a theoretical explanation: under CFG, raising the class-conditional density to the power $(1+\lambda)$ disproportionately collapses modes of high-entropy (balanced multi-modal) classes, making them more sensitive to the guidance scale and more likely to benefit from adaptive rather than fixed guidance. Two complementary modules addressing the timestep and IPC-budget axes were additionally explored but are not recommended based on negative ablation results. \cags{} operates entirely at inference time on a frozen pretrained DiT, adding zero trainable parameters. Experiments on ImageNet-100, ImageNette, and ImageWoof across ConvNet-6, ResNet-18, and ResNetAP-10 demonstrate that \cags{} yields consistent and substantial improvements on the 100-class subset: $+7.6\%$ relative at IPC=10 and $+14.8\%$ at IPC=50, with gains scaling on deeper architectures ($+17.0\%$ on ResNet-18, $+13.4\%$ on ResNetAP-10). On the 10-class subset, \cags{} provides a large gain at IPC=50 ($+9.0\%$) but is neutral at IPC=10 and \emph{harmful} at IPC=1 ($-12.8\%$). A cross-dataset analysis reveals a scale-dependent trade-off: entropy-only guidance outperforms the multi-factor metric on small-scale datasets ($\leq$10 classes) by $1.4$--$1.6\%$, while the multi-factor \cags{} metric wins on medium-to-large-scale datasets ($\geq$44 classes) by $0.65$--$0.94\%$. A semantic subset analysis on ImageNet-100 partitions confirms this trade-off is driven by class count, not semantic composition. Scaling to 200 classes amplifies the \cags{} advantage from $+2.50\%$ to $+3.57\%$ absolute ($+19.2\%$ relative), confirming that per-class adaptation becomes increasingly beneficial as complexity variation grows. A fair comparison with D4M under the same protocol shows that real-image initialization ($20.21\%$) underperforms even unguided generation ($21.63\%$), confirming that the strength of diffusion-based distillation lies in the generative prior rather than preserved real-image features.
\end{abstract}